% !TeX TXS-program:compile = txs:///lualatex | txs:///view
\documentclass{article}
\usepackage[spanish]{babel}
\usepackage{amsmath}
\usepackage{amssymb}
\usepackage{lipsum}

\input{colores} % Apunta al archivo.
\input{dict_ecuaciones}
\input{cajas}
\input{quimica}

\newtheorem{definicion}{Definición}
\newtheorem{obs}{Observación}[section]
\title{Tomando el control de \LaTeX}
\author{Oromion}
\date{\today}

\begin{document}

\maketitle

\textcolor{myGreen}{un texto en color verde}

\lipsum[1]

\cuadratica{}

\lipsum[1]

\cuadratica{}

\Fourier[eq]

\section{Sección 1}

\lipsum[1]
\begin{definicion}
\lipsum[1]
\end{definicion}

\begin{obs}
Esta es una primera
\end{obs}

\begin{definicion}
	\lipsum[1]
\end{definicion}


\begin{obs}
Segunda observación

\end{obs}

\section{Sección 2}

\begin{obs}
Observación 2.1
\end{obs}
\end{document}

Veremos ahora los comandos y su redefinición.